\section*{Overview}

Most contest geometry starts with one predicate: given three points, do they make a left turn, a right turn, or stay
collinear? That question is answered by the cross product, and once that primitive is solid, many other routines become
simple compositions of it.

\section*{When to Use It}

Use orientation and cross products when:

\begin{itemize}
  \item you need to compare turns without floating-point angles,
  \item the task asks about segment intersection, convexity, polygon area, or hull building,
  \item integer coordinates are available and exact predicates matter.
\end{itemize}

\section*{Core Idea}

For vectors \(\vec{u} = (x_1, y_1)\) and \(\vec{v} = (x_2, y_2)\), the 2D cross product is
\[
\vec{u} \times \vec{v} = x_1 y_2 - y_1 x_2.
\]

For three points \(a, b, c\), compute
\[
(b - a) \times (c - a).
\]
Its sign tells the turn:
\begin{itemize}
  \item positive: left turn,
  \item negative: right turn,
  \item zero: collinear.
\end{itemize}

\section*{Key Insight}

The cross product is both a turn test and a signed-area test. Its magnitude is twice the signed area of triangle
\(abc\). That one fact explains why it appears in hulls, polygon area, intersection predicates, and half-plane tests.

\section*{Operations / Main Technique}

\begin{itemize}
  \item \textbf{orientation}: left / right / collinear.
  \item \textbf{on-segment test}: after confirming collinearity, check whether a point lies inside the bounding box.
  \item \textbf{polygon area}: sum edge cross products.
  \item \textbf{turn filtering}: keep only left turns when building a convex hull.
\end{itemize}

\paragraph{Worked example.}
If \(a = (0,0)\), \(b = (3,0)\), and \(c = (1,2)\), then
\[
(b-a) \times (c-a) = 3 \cdot 2 - 0 \cdot 1 = 6 > 0,
\]
so \(a \to b \to c\) is a left turn.

\section*{Correctness Intuition}

The sign of the cross product tells which side of the directed line \(a \to b\) the point \(c\) lies on. Everything
else is built on top of that side test. Since the computation uses only integer arithmetic, the predicate is exact as
long as overflow is avoided.

\section*{Complexity Analysis}

Each primitive here is \(O(1)\) time and \(O(1)\) memory.

\section*{Implementation}

The reference code provides:

\begin{itemize}
  \item a point type on 64-bit integers,
  \item \texttt{cross(a, b, c)} for orientation,
  \item \texttt{orient} returning \(-1, 0, 1\),
  \item \texttt{on\_segment} for boundary checks.
\end{itemize}

\section*{Common Pitfalls}

\begin{itemize}
  \item Using \texttt{int} when coordinates can make the cross product overflow.
  \item Forgetting that a zero cross product only means collinear, not necessarily on the segment.
  \item Mixing clockwise and counterclockwise conventions between notes and code.
  \item Reaching for \(\texttt{atan2}\) when a sign test would be simpler and more reliable.
\end{itemize}

\section*{Variants / Extensions}

\begin{itemize}
  \item Dot product for projection and angle-related tests.
  \item Shoelace formula for polygon area.
  \item Orientation as the core primitive for segment intersection and convex hull.
\end{itemize}

\section*{Practice Problems}

\begin{itemize}
  \item Determine whether two segments intersect.
  \item Compute polygon area from its vertices.
  \item Filter turns while constructing a convex hull.
  \item Check whether a polygon is strictly convex.
\end{itemize}

\section*{References}

\begin{itemize}
  \item \href{https://cp-algorithms.com/geometry/basic-geometry.html}{cp-algorithms: Basic Geometry}.
  \item \href{https://usaco.guide/plat/geo-prims}{USACO Guide: Geometry Primitives}.
  \item \href{https://codeforces.com/catalog}{Codeforces Catalog}.
\end{itemize}
